%% SCRAP: experiments/02-experiments/L8_WIRING_CODE_REVIEW_GUIDE
%% SOURCE: docs/working/experiments/02-experiments/L8_WIRING_CODE_REVIEW_GUIDE.md
%% STATUS: HISTORICAL
%% FITS: experiments/ch-l8
%% EDITORIAL: lifted — prose rewritten to press voice

\section{L8 Control-Wiring Code Review}
\label{sec:l8-wiring-code-review}

This section documents the surgical code changes that completed the final
``last mile'' connection between the L8 Jacquard mode selector and the four
adaptive feedback loops it governs (L2, L3, L5, L6). The changes span three
source files and approximately 55 lines of new conditional logic.

\subsection{Summary}

All 731 implemented tests passed on a clean \texttt{make fastest} build with
zero warnings and zero errors. The four execution-path gates were verified
individually and in combination.

\subsection{Files Modified}

\subsubsection{\texttt{src/physics\_metadata.c} — L3 Linear Decay Gate}

The gate is inserted in \texttt{physics\_metadata\_apply\_linear\_decay()} after
the existing null, frozen-word, and minimum-interval checks, placing it at the
optimal early-exit position:

\begin{lstlisting}[language=C]
/* L8 GATE: Check if L3 (linear decay) is enabled by Jacquard mode selector */
if (vm->ssm_config) {
    ssm_config_t *config = (ssm_config_t*)vm->ssm_config;
    if (!config->L3_linear_decay) {
        return;  /* L3 disabled by L8 mode selector */
    }
}
\end{lstlisting}

When L3 is disabled, words retain execution heat without temporal decay, freezing
the thermal state until the mode selector re-enables the loop.

\subsubsection{\texttt{src/inference\_engine.c} — L5 and L6 Inference Gates}

Both gates appear in \texttt{inference\_engine\_run()}, one at Phase~2C
(window-width inference) and one at Phase~2D (decay-slope inference). Each gate
preserves the previous output value when the loop is disabled, providing
\emph{sticky configuration} semantics: the last computed value remains in effect
rather than resetting to a default.

\begin{lstlisting}[language=C]
/* L5: Window Width Inference gate */
uint32_t inferred_width = outputs->adaptive_window_width; /* keep previous */
if (inputs->vm && inputs->vm->ssm_config) {
    ssm_config_t *config = (ssm_config_t*)inputs->vm->ssm_config;
    if (config->L5_window_inference) {
        inferred_width = find_variance_inflection(
            trajectory, traj_len, current_variance);
    }
} else {
    inferred_width = find_variance_inflection(  /* legacy mode */
        trajectory, traj_len, current_variance);
}
\end{lstlisting}

The L6 gate is structurally identical, gating \texttt{infer\_decay\_slope\_q48()}.
When either inference loop is off, disabling it avoids Levene's test
($\sim$5{,}000~cycles saved) or exponential regression ($\sim$3{,}000~cycles saved)
on every heartbeat tick.

\subsubsection{\texttt{src/vm.c} — L2 Rolling Window Gate}

The L2 gate sits in the hot path of \texttt{execute\_colon\_word()}, adding a
null check and one boolean read ($\sim$2~CPU cycles) per word execution:

\begin{lstlisting}[language=C]
if (vm->ssm_config) {
    ssm_config_t *config = (ssm_config_t*)vm->ssm_config;
    if (config->L2_rolling_window) {
        rolling_window_record_execution(&vm->rolling_window, word_id);
    }
} else {
    rolling_window_record_execution(&vm->rolling_window, word_id);
}
\end{lstlisting}

Call-site gating (Option~A) was chosen over an internal gate inside
\texttt{rolling\_window\_record\_execution()} to preserve separation of concerns.

\subsection{L8 Signal Path}

The complete control chain from measurement to enforcement is:

\begin{enumerate}
  \item \texttt{vm\_heartbeat\_update\_l8()} computes entropy, CV, temporal
        decay, and stability score each heartbeat cycle.
  \item \texttt{ssm\_l8\_update()} evaluates the four metrics against thresholds
        and selects a 4-bit mode (C0--C15) subject to 5-tick hysteresis.
  \item \texttt{ssm\_apply\_mode()} decodes the mode bits into the L2/L3/L5/L6
        boolean flags in \texttt{vm->ssm\_config}.
  \item The execution-path gates (this implementation) read those flags at
        runtime.
\end{enumerate}

\subsection{Performance Overhead}

\begin{center}
\begin{tabular}{lrr}
\toprule
Loop & Gate overhead & Saving when disabled \\
\midrule
L2 (rolling window) & $\sim$2 cycles & $\sim$10 cycles \\
L3 (linear decay)   & $\sim$2 cycles & $\sim$50 cycles \\
L5 (window inference) & $\sim$2 cycles & $\sim$5{,}000 cycles \\
L6 (decay inference)  & $\sim$2 cycles & $\sim$3{,}000 cycles \\
\bottomrule
\end{tabular}
\end{center}

Net impact across all four gates is less than 1\% overhead; the savings when
loops are suppressed are substantial.

\subsection{Default Behaviour}

The VM initialises with mode C0 (all loops off). Within approximately 25~ms the
L8 selector adapts to the running workload and enables the appropriate loop
subset. The test suite passes in mode~C0, confirming that the gates do not break
any existing behaviour.
